\chapter{Linear Diophantine Equations}\label{chapter:linear-diophantine-equations}

In this section we study \emph{linear Diophantine equations}. These are equations of the form
\begin{equation*}
  a_0x_0+\cdots +a_kx_k=b,
\end{equation*}
with integer coefficients $a_0,\ldots,a_k$. The objective of a linear Diophantine equation such as the one above is to find integers $x_0,\ldots,x_k$ for which the equation holds, to study conditions under which the equation has integer solutions, and ultimately to characterize all solutions.

\section{Greatest Common Divisors}

Some preorders $(P,\leq)$ satisfy the property that for any two elements $a,b\in P$ there is a maximal element below both of them. This property can be formulated succinctly as follows: An element $d$ is a maximal element below $a$ and $b$ if it satisfies the logical equivalence
\begin{equation*}
  (x \leq a) \land (x \leq b) \Leftrightarrow (x \leq d)
\end{equation*}
for every element $x\in P$. If $d$ satisfies this property, we also say that $d$ is the \emph{meet} of $a$ and $b$. It is common to write $a\wedge b$ for the meet of $a$ and $b$, if $P$ is a poset in which any two elements have a meet.

The definition of meets needs some explanation. It uses the \emph{conjunction} ($\land$) and it uses \emph{logical equivalence} ($\Leftrightarrow$). The conjunction $p\land q$ of two propositions $p$ and $q$ is the proposition that both $p$ and $q$ are true. The logical equivalence $p \Leftrightarrow q$ is the proposition that $p$ is true if and only if $q$ is true. In other words, it is the conjunction of the proposition that $p$ implies $q$ and the proposition that $q$ implies $p$.

If $d$ is the meet of $a$ and $b$, then the logical equivalence gives two implications back and forth:
\begin{enumerate}
\item Any element $x$ that is below both $a$ and $b$, must be below $d$.
\item Any element $x$ that is below $d$, must be below both $a$ and $b$. In particular, $d$ itself satisfies $d\leq d$, so it follows that $d\leq a$ and $d\leq b$. That is, $d$ itself is below both $a$ and $b$, and by the first property it is the maximal such element.
\end{enumerate}

We wouldn't have this discussion if it didn't apply for us. Indeed, we will show that the preorder of integers with divisibility has meets. In the case of the integers, we will call such meets \emph{greatest common divisors}.

\begin{definition}
  Consider two integers $a$ and $b$. An integer $d$ is said to be a \define{greatest common divisor} of $a$ and $b$ if it satisfies the following logical equivalence
  \begin{equation*}
    (x\mid a)\land(x\mid b) \Leftrightarrow (x\mid d)
  \end{equation*}
  for every integer $x$.
\end{definition}

Thus, a greatest common divisor of $a$ and $b$ is an integer $d$ that divides both $a$ and $b$, and any integer that divides both $a$ and $b$ is a divisor of $d$. In other words, the signifier "greatest" refers to the divisibility ordering of the integers.

\begin{example}
  The greatest common divisor of $0$ and $0$ is $0$. To see this, simply note that the proposition
  \begin{equation*}
    (x\mid 0)\land(x\mid 0) \Leftrightarrow (x\mid 0)
  \end{equation*}
  is true because any proposition is logically equivalent to the conjunction with itself. Here we see why it was relevant to note that the word "greatest" refers to the divisibility ordering: Any integer is a divisor of $0$, but none of them are the greatest in the standard ordering $\leq$.
\end{example}

\begin{theorem}\label{theorem:definition-gcd}
  Any two integers $a$ and $b$ have a greatest common divisor, for which we write $\gcd(a,b)$. Furthermore, there are integers $k$ and $l$ satisfying \emph{B\'ezout's identity}
  \begin{equation*}
    ax+by = \gcd(a,b).
  \end{equation*}
\end{theorem}

\begin{proof}
  Consider two integers $a$ and $b$, and let $I:=(a)$ and $J:=(b)$ be the principal ideals generated by $a$ and $b$. Now consider the set
  \begin{equation*}
    I+J := \{x+y \mid x\in I,\,y\in J\}
  \end{equation*}
  of all sums of elements of $I$ and elements of $J$. We claim that the set $I+J$ is again an ideal. It contains $0$, because $0=0+0$ is a sum of elements in $I$ and $J$. Furthermore, to see that $I+J$ is closed under linear combinations of its elements, consider two elements $u,v\in I+J$. If $u=x_u+y_u$ and $v=x_v+y_v$, then we can write any linear combination $ku+lv$ of $u$ and $v$ in the form
  \begin{equation*}
    ku+lv = (kx_u + lx_v) + (ky_u + ly_v),
  \end{equation*}
  which is the sum of an element in $I$ and an element in $J$. Thus, it follows that $I+J$ is an ideal.

  Now let $d$ be the unique natural number such that $I+J=(d)$, which exists by \zcref{theorem:integers-principal-ideal-domain}. We claim that $d$ is the greatest common divisor of $a$ and $b$. Notice that $a=a+0\in I+J$, so $a$ must be a multiple of $d$. Similarly, we have $b=0+b\in I+J$, so $b$ must be a multiple of $d$. In other words, $d$ is a common divisor of $a$ and $b$. Furthermore, if $e$ is any common divisor of $a$ and $b$, then $e$ is a divisor of any element in $I+J$. In particular, we must have $e\mid d$, allowing us to conclude that $d$ is the greatest common divisor.

  The last claim follows, because $d$ is an element of $I+J$, which consists elements of the form $ax+by$.
\end{proof}

A very important condition on pairs of integers $a$ and $b$, which occurs as an assumption in many lemmas, propositions, and theorems, is the condition that $\gcd(a,b)=1$. In other words, that any common divisor of $a$ and $b$ must be $1$. Such numbers are called \emph{relatively prime} or \emph{coprime}.

\begin{definition}
  Two integers $a$ and $b$ are said to be \define{relatively prime} or \define{coprime} if
  \begin{equation*}
    \gcd(a,b)=1.
  \end{equation*}
\end{definition}

The following proposition is a must-know fact about the integers.

\begin{proposition}\label{proposition:coprime-divides-product}
  Consider two relatively prime integers $a$ and $b$, and an arbitrary integer $c$. Then we have
  \begin{equation*}
    a\mid bc \Leftrightarrow a\mid c.
  \end{equation*}
\end{proposition}

\begin{proof}
  The fact that $a\mid c$ implies $a\mid bc$ is obvious, since if $c$ is a multiple of $a$, then so is $bc$. This establishes the converse direction of the logical equivalence.

  For the forward direction, assume that $a\mid bc$. Since we have assumed that $a$ and $b$ are relatively prime, there is a solution to the equation
  \begin{equation*}
    ax+by=1.
  \end{equation*}
  Consequently, there is a solution to the equation $cax+cby=c$. However, both $cax$ and $cby$ are divisible by $a$, so $c$ must be divisible by $a$.
\end{proof}

\section{Euclid's Algorithm}

\zcref{theorem:definition-gcd} shows that any two integers $a$ and $b$ have a greatest common divisor $\gcd(a,b)$, and that there are integers $k$ and $l$ such that
\begin{equation*}
  ka+lb=\gcd(a,b),
\end{equation*}
but it doesn't reveal much about how the integers $k$ and $l$ can be found, or even how to find the greatest common divisor in the first place. This problem is solved by \emph{Euclid's Algorithm}, which provides an efficient way to compute the greatest common divisor of two integers and to find integers $k$ and $l$ establishing the greatest common divisor as a linear combination. The algorithm relies on the following proposition.

\begin{proposition}\label{proposition:gcd-euclidean-division}
  Suppose $a$ and $b$ are two integers and $b=qa+r$. Then we have
  \begin{equation*}
    \gcd(b,a)=\gcd(a,r).
  \end{equation*}
\end{proposition}

\begin{proof}
  Suppose that $d$ is a divisor of $a$. Then we claim that $d\mid b$ if and only if $d\mid r$. Indeed, since $b=qa +r$ it follows that $b$ is a linear combination of $a$ and $r$, and also that $r$ is a linear combination of $a$ and $b$. Thus, any divisor of $a$ and $b$ is also a divisor of $r$, and any divisor of $a$ and $r$ is also a divisor of $b$. In other words, we have
  \begin{equation*}
    d\mid b \Leftrightarrow d\mid r.
  \end{equation*}
  It follows that $d$ is a common divisor of $a$ and $b$ if and only if $d$ is a common divisor of $a$ and $r$. This implies that their greatest common divisors coincide.
\end{proof}

As an immediate corollary, we obtain:

\begin{corollary}\label{corollary:gcd-euclidean-division}
  If $b$ is of the form $ka+1$, then $a$ and $b$ are always relatively prime.
\end{corollary}

Euclid's algorithm works by repeatedly using \zcref{proposition:gcd-euclidean-division}. For any pair $(a,b)$ of integers, if $b\neq 0$ then there is a unique pair $(q,r)$ such that $0\leq r<|b_0|$ and such that the equality
\begin{equation*}
  a=qb+r
\end{equation*}
holds. This determines a new pair $(a',b'):=(b,r)$ and $\gcd(a,b)=\gcd(b,r)$. Note that $\gcd(b,r)$ is somewhat easier to compute, since the numbers involved are smaller. Repeating this procedure, we obtain a list of equalities
\begin{align*}
  a_1 & =q_1b_1+r_1, \\
  a_2 & =q_2b_2+r_2, \\
      & ~\,\vdots \\
  a_N & = q_Nb_N+r_N,
\end{align*}
in which $a_1=a$, $b_1=b$, $a_{i+1}=b_i$, $b_{i+1}=r_i$, and $0\leq r_i<|b_i|$. This procedure results in a decreasing sequence $r_1>r_2>\cdots>r_N$ of nonnegative integers, so it terminates when $r_N=0$. Indeed, the Euclidean division theorem cannot be applied to the pair $(b_N,r_N)$ any further when $r_N=0$. Since \zcref{proposition:gcd-euclidean-division} implies that
\begin{equation*}
  \gcd(a,b)=\gcd(b_1,r_1)=\gcd(b_2,r_2)=\cdots=\gcd(b_N,r_N),
\end{equation*}
we see that $b_N$ is the greatest common divisor of $a$ and $b$.

\begin{example}
  We will illustrate Euclid's algorithm by an example. Suppose we want to calculate the greatest common divisor of $578$ and $732$. Then we write $732$ in the form $q\cdot 578+r$; that is, we write
  \begin{equation*}
    732 = 1 \cdot 578 + 154.
  \end{equation*}
  By \zcref{proposition:gcd-euclidean-division}, it follows that $\gcd(732,578)=\gcd(578,154)$. We now proceed in the same manner, by writing
  \begin{equation*}
    578 = 3\cdot 154 + 116, 
  \end{equation*}
  and we observe that $\gcd(732,578)=\gcd(154,116)$. We proceed with this process until we find that $\gcd(732,578)=\gcd(d,0)$ for some number $d$. Once we reach this stage, we conclude that $\gcd(732,578)=d$. Putting words to action:
  \begin{align*}
    154 & = 1\cdot 116 + 38 & \gcd(732,578) & = \gcd(116,38), \\
    116 & = 3 \cdot 38 + 2 & \gcd(732,578) & = \gcd(38,2), \\
    38 & = 19 \cdot 2 + 0 & \gcd(732,578) & = \gcd(2,0).
  \end{align*}
  Thus, we conclude that $\gcd(732,578)=2$.  
\end{example}

A few further remarks about the sequence of equations
\begin{equation*}
  a_i=q_ib_i+r_i
\end{equation*}
produced by Euclid's algorithm are in order. Since each $r_i$ can be written as a linear combination of $a_i$ and $b_i$, it follows by induction that each $r_i$ can be written as a linear combination of $a$ and $b$. In other words, for each $1\leq i\leq N$, there are integers $k_i$ and $l_i$ such that
\begin{equation*}
  r_i=k_ia+l_ib.
\end{equation*}
These linear coefficients can be determined recursively, by what is called the \emph{extended Euclidean algorithm}. Notice that $k_1=1$ and $l_1=q_1$. We also set $k_0:=0$ and $l_0:=1$. Then $k_{i+1}$ and $l_{i+1}$ are recursively determined by the rule:
\begin{align*}
  k_{i+1} & = k_{i-1}-q_{i+1}k_i, \\
  l_{i+1} & = l_{i-1}-q_{i+1}l_i.
\end{align*}
To see this, simply note that
\begin{equation*}
  r_{i+1}=a_{i+1}-q_{i+1}b_{i+1}=r_{i-1}-q_{i+1}r_i.
\end{equation*}

\begin{example}
  In the previous example, Euclid's algorithm has given us that
  \[\gcd(732,578)=2.\]
  We will now express $2$ as a linear combination of $732$ and $578$. We start with the first equation, $732=1\cdot 578+154$ to write
  \begin{align*}
    154 & = 732 - 1\cdot 578.
  \end{align*}
  Next, we use the equation $578=3\cdot 154+116$ to write
  \begin{equation*}
    116 = 578-3\cdot 154 = 578 - 3\cdot (732 - 1\cdot 578) = 4\cdot 578 - 3\cdot 732.
  \end{equation*}
  Continuing this way, we find that
  \begin{equation*}
    38=154-1\cdot 116 = (732 - 1\cdot 578) - 1\cdot (4\cdot 578 - 3\cdot 732) = 4\cdot 732 - 5\cdot 578
  \end{equation*}
  and
  \begin{equation*}
    2 = 116-3\cdot 38 = (4\cdot 578 - 3\cdot 732) - 3\cdot (4\cdot 732 - 5\cdot 578) = 19\cdot 578 - 15\cdot 732.
  \end{equation*}
  In the second to last equation, we find that $2$ is a linear combination of $38$ and $116$; that is,
\begin{align*}
  2 & = 116 - 3\cdot 38.
\end{align*}
However, by the equation before that we see that $38$ itself is a linear combination of $154$ and $116$. This can be used to express $2$ as a linear combination of $154$ and $116$; that is,
\begin{align*}
  2 & = 116 - 3\cdot (154 - 116) & \Rightarrow & & 2 & = 4\cdot 116 - 3\cdot 154.
\intertext{Furthermore, $116$ is a linear combination of $154$ and $578$, giving us that}
  2 & = 4\cdot (578-3\cdot 154) - 3\cdot 154 & \Rightarrow & & 2 & = 4\cdot 578 - 15\cdot 154.
\intertext{In the final step, we use that $154$ is a linear combination of $732$ and $578$ to find}
  2 & = 4\cdot 578 - 15\cdot (732 - 578) & \Rightarrow & & 2 & = - 15\cdot 732 + 19\cdot 578.
\end{align*}
\end{example}

\begin{theorem}\label{theorem:correspondence-theorem-quotients-euclidean-algorithm}
  Every nonempty list $(q_0,\ldots,q_N)$ of positive integers uniquely determines a pair $(a,b)$ of relatively prime positive integers such that $a>b$. In other words, there is a bijection
  \begin{equation*}
    \{(q_0,\ldots,q_N)\mid q_i\geq 1\text{ for all }0\leq i \leq N\} \cong \{(a,b)\mid \gcd(a,b)=1\text{ and }a>b>0\}.
  \end{equation*}
\end{theorem}

\section{Linear Diophantine Equations in Multiple Variables}

In this section we turn our attention to linear Diophantine equations in two variables, that is, integer equations of the form
\begin{equation*}
  ax + by = c.
\end{equation*}
We begin by applying the theory of ideals to establish that this equation has a solution if and only if $\gcd(a,b)\mid c$.

\begin{theorem}\label{theorem:solve-linear-diophantine-equation}
  Consider integers $a$, $b$, and $c$, and let $d=\gcd(a,b)$. The linear Diophantine equation
  \begin{equation*}
    ax + by = c
  \end{equation*}
  is solvable with integers $x$ and $y$ if and only if $d\mid c$.
\end{theorem}

\begin{proof}
  The first part of the claim is equivalent to the claim that we have an equality of ideals
  \begin{equation*}
    (a)+(b) = (\gcd(a,b)).
  \end{equation*}
  We recall that the ideal $I+J$ consists of sums $x+y$ of $x\in I$ and $y\in J$. Furthermore, the ideals $(a)$ and $(b)$ consist of multiples of $a$ and $b$, respectively. Thus the ideal $(a)+(b)$ consists of all the linear combinations of $a$ and $b$; in other words, all the possible integers $c$ for which there is a solution to the equation $ax+by=c$.

  In \zcref{theorem:definition-gcd} we defined $\gcd(a,b)$ to be the unique natural number $d$ such that $(a)+(b)=(d)$. Such a natural number exists by \zcref{theorem:integers-principal-ideal-domain}. Now we observe that the equality
  \begin{equation*}
    (a)+(b)=(\gcd(a,b))
  \end{equation*}
  tells us exactly that the set of integers $c$ for which there exists a solution of the equation $ax+by=c$ has the same elements as the set of multiples of $\gcd(a,b)$. In other words, a solution to $ax+by=c$ exists if and only if $c$ is a multiple of the greatest common divisor of $a$ and $b$.
\end{proof}

As an immediate corollary of the previous theorem, we note that:

\begin{corollary}
  If $a$ and $b$ are relatively prime integers, then the linear Diophantine equation
  \begin{equation*}
    ax + by = c
  \end{equation*}
  is solvable for any $c$.
\end{corollary}

Now that we have established a necessary and sufficient condition for its solvability, it remains to describe a way of finding all solutions, if there are any. Notice that we have already collected quite a bit of useful information. Euclid's algorithm allows find a solution to the equation
\begin{equation*}
  ax + by = d,
\end{equation*}
where $d=\gcd(a,b)$. Thus, if $c=kd$, then we can solve
\begin{equation*}
  ax + by = c
\end{equation*}
by first finding a solution $ax_0+by_0 = d$ by Euclid's algorithm. Then we find that
\begin{equation*}
  a(kx_0)+b(ky_0)=kd = c.
\end{equation*}
In other words, we can use Euclid's algorithm to find a solution to the equation $ax+by=c$ if it has any.

\begin{proposition}
  Consider integers $a$, $b$, and $c$, and assume that $a$ and $b$ are relatively prime. If $ax_0 + by_0=c$ is a solution to the Diophantine equation
  \begin{equation*}
    ax + by = c,
  \end{equation*}
  then every solution is of the form
  \begin{equation*}
    x = x_0 + bk,\qquad\text{and}\qquad y = y_0-ak.
  \end{equation*}
\end{proposition}

\begin{proof}
  Suppose that $ax+by=c$ is another solution. Then it follows that
  \begin{equation*}
    a(x-x_0)+b(y-y_0)=0.
  \end{equation*}
  In other words, we have $a(x-x_0)=-b(y-y_0)$. Thus we see that $b(y-y_0)$ is divisible by $a$. Since $a$ and $b$ are relatively prime, this implies that $a\mid y-y_0$. Similarly, $a(x-x_0)$ is divisible by $b$, and therefore it follows that $b\mid x-x_0$. Furthermore, if $ka=y-y_0$ and $lb=x-x_0$. Then the equation
  \begin{equation*}
    lab - kab = 0
  \end{equation*}
  implies that $k=l$. In other words, $x=x_0+bk$ and $y=y_0-ak$.
\end{proof}

\begin{theorem}\label{theorem:set-of-solutions-binary-linear-diophantine-equation}
  Suppose that $ax_0,by_0=c$ is a solution to the Diophantine equation
  \begin{equation*}
    ax + by = c.
  \end{equation*}
  Then every solution is of the form
  \begin{equation*}
    x = x_0 + \frac{b}{d}k,\qquad\text{and}\qquad y = y_0-\frac{a}{d}k.
  \end{equation*}
\end{theorem}

\begin{proof}
  This theorem follows from the previous proposition. If $d=\gcd(a,b)$, then we find that
  \begin{equation*}
    \frac{a}{d}x+\frac{b}{d}y=\frac{c}{d}.
  \end{equation*}
  The integers $\frac{a}{d}$ and $\frac{b}{d}$ are relatively prime, so we find that $x=x_0+\frac{b}{d}k$ and $y=y_0-\frac{a}{d}k$ for some integer $k$.
\end{proof}

\section{Āryabhaṭa's Computations of the Positions of Planets}

% \section{The Kuṭṭaka Method}

\exercises

\begin{starter}
  \exercise Suppose that $d$ is a common divisor of two nonzero integers $a$ and $b$.
  \begin{subexenum}
  \item Show that the greatest common divisor of $a/d$ and $b/d$ is $\gcd(a,b)/d$.
  \item Show that $d$ is a greatest common divisor of $a$ and $b$ if and only if $a/d$ and $b/d$ are relatively prime.  
  \end{subexenum}
  \exercise Show that $a\mid bc$ if and only if $a/\gcd(a,b) \mid c$.
  \exercise Show that if $\gcd(a,b)=1$ and both $a\mid c$ and $b\mid c$, then $ab\mid c$.  
  \exercise Show that $\gcd(a^n,b^n)=\gcd(a,b)^n$.
  \exercise Show that if $d\mid a + b$, then $\gcd(d,a)\mid \gcd(a,b)$ and $\gcd(d,b)\mid\gcd(a,b)$. Conclude that if $a$ and $b$ are relatively prime, then $d$ is relatively prime to both $a$ and $b$.
  \exercise Consider three integers $a$, $b$, and $c$. Show that the following are equivalent:
  \begin{enumerate}
  \item The integers $a$ and $b$ are both relatively prime to $c$.
  \item The integer $ab$ is relatively prime to $c$.
  \end{enumerate}
\end{starter}

\begin{routine}
  \exercise Show that $a\mid b$ implies $ac\mid bc$. Furthermore, show that if $c$ is nonzero, then $ac\mid bc$ implies $a\mid b$. Conclude that $a^n\mid b^n$ if and only if $a\mid b$ for any nonzero natural number $n$.
  \exercise
  \begin{subexenum}
  \item Show that any three consecutive Fibonacci numbers are pairwise relatively prime.
  \item Find all $n$ for which the five consecutive Fibonacci numbers $F_n,...,F_{n+4}$ are pairwise relatively prime.
  \end{subexenum}
  \exercise Show that $n!+1$ and $(n+1)!+1$ are relatively prime for all nonzero $n$.
\end{routine}

\begin{challenge}
  \exercise \label{exercise:div-increment-a-p-squares}
  \begin{subexenum}
  \item Show that if $a$ and $b$ are both relatively prime to $24$, then $a^2-b^2$ is divisible by $24$.
  \item A \emph{generalized pentagonal number} is a number of the form
    \begin{equation*}
      \frac{3k^2\pm k}{2}.
    \end{equation*}
    Show that if $a$ is the $n$\supth{} integer relatively prime to $24$, then the number
    \begin{equation*}
      \frac{a^2-1}{24}.
    \end{equation*}
    is the $n$\supth{} generalized pentagonal number.
  \item Let $0\leq a\leq b\leq c$ be integers such that their squares form an arithmetic progression; that is, the numbers $a^2$, $b^2$, and $c^2$ listed in increasing order have a common increment:
  \begin{equation*}
    \Delta:=b^2-a^2=c^2-b^2.
  \end{equation*}
  Show that $\Delta$ is a multiple of $24$ \cite{fibonacci1987}.
  \item Find five examples of triples $(a,b,c)$ as described in part (a), and explain how you found them.
  \end{subexenum}
  \exercise
  \begin{subexenum}
  \item Consider two positive integers $a$ and $n$ such that $\gcd(a,n)=1$. Prove that there is a natural number $m$ such that
    \begin{equation*}
      n\,\Bigm|\, \sum_{k=0}^m a^k.
    \end{equation*}
    Hint: Apply the pigeonhole principle.
  \item Find the least positive integer $m$ such that
    \begin{equation*}
      7\,\Bigm|\, \sum_{k=0}^{m-1} 10^k = \underbrace{11\cdots1}_{m}.
    \end{equation*}
  \end{subexenum}
  \exercise Show that $2^m-1$ and $2^n+1$ are relatively prime, provided that $m$ is odd.
  \exercise Show that
  \begin{equation*}
    \binom{a+k}{k}\Bigm| \lcm(a,a+1,\ldots,a+k).
  \end{equation*}
  \exercise (Erd\H{o}s) Show that any set of $n+1$ integers chosen from an interval $a,\ldots,a+2n-1$ must contain a pair of relatively prime integers.
  \exercise An interval $a,\ldots,a+n-1$ of consecutive positive integers is said to be \emph{stapled} if it does not contain any integer that is relatively prime to all the others \cite{evans1969}. The sequence A090318 of the OEIS lists for each $n$ the first positive integer $a$ such that the interval $a,\ldots,a+n-1$ is stapled, or $0$ if no such number exists. This sequence starts to show nonzero values from $n=17$ onwards, with the value $2184$ at $n=17$.
  \begin{subexenum}
  \item Check that the interval $2184,\ldots, 2200$ is stapled.
  \item Show that in every interval of ten consecutive digits there is an integer that is relatively prime to all the others. In other words, no interval of ten consecutive digits is stapled.
  \end{subexenum}
\end{challenge}

%%% Local Variables:
%%% mode: latex
%%% TeX-master: "elementary-number-theory"
%%% TeX-engine: luatex
%%% End:
